\section{Thu hẹp vector ngữ cảnh, và xem chỗ nào hỏng trước}
\label{sec:ch05-bop-vector}

\index{vector ngữ cảnh!đo bằng cách thu hẹp}%
\tn{Vector ngữ cảnh}{context vector} có bề rộng cố định, và ai cũng nói đó là
chỗ thắt. Câu ấy dễ nói và khó kiểm, vì trên một mô hình thật thì bề rộng
vector dính chặt với số tham số, với sức chứa của encoder, với sức chứa của
decoder, và một mô hình hẹp hơn thì kém hơn ở mười lý do cùng lúc.

Cách tách chúng ra là giữ nguyên tất cả rồi chỉ đổi một thứ. Cùng corpus, cùng
optimizer, cùng seed, cùng số lần cập nhật; chỉ \(d_h\) đổi, và decoder nhận
đúng \(2 d_h\) số thực dù \tn{câu nguồn}{source sentence} dài bao nhiêu.

\begin{measured}{độ chính xác theo bề rộng vector ngữ cảnh}
\begin{minted}{text}
d_hidden  context  params   loss     exact    short    long
4         8        3641     1.5756   0.0000   0.0000   0.0000
8         16       5229     1.1748   0.0000   0.0000   0.0000
16        32       9173     0.8733   0.0033   0.0068   0.0000
32        64       20133    0.5584   0.1433   0.2905   0.0000
64        128      54341    0.1429   0.5200   0.9797   0.0724
128       256      171909   0.0251   0.8533   1.0000   0.7105
\end{minted}

6000 cặp huấn luyện, 300 cặp test, 14 epoch, câu nguồn đảo, giải mã tham lam.
Một seed mỗi dòng.

Đo bằng \code{experiments/ch05\_bottleneck.py} tại tag \code{ch05}.
\end{measured}

\index{vector ngữ cảnh!hỏng không đều theo độ dài}%
Cột \code{exact} đi từ 0 lên 0.8533 và đó là chuyện dự đoán được. Hai cột cuối
mới là chỗ đáng dừng lại.

Vector hẹp \emph{không} hỏng đều. Ở \(d_h = 64\), tức 128 số thực, câu ngắn đã
đạt 0.9797 trong khi câu dài còn 0.0724. Cùng một mô hình, cùng một lần huấn
luyện, chênh nhau hơn chín mươi điểm, và thứ duy nhất phân biệt hai nhóm là câu
nguồn dài bao nhiêu token. Ở \(d_h = 32\) thì câu ngắn được 0.2905 còn câu dài
tròn 0 không sót câu nào.

Đó là hình dạng của câu \enquote{câu phải nhét vừa vào vector}, khi đo thay vì
khi phát biểu. Một vector rộng \(2 d_h\) số thực chứa được câu ngắn trước, và
ngưỡng mà nó bắt đầu chứa nổi câu dài nằm cao hơn hẳn. Giữa hai ngưỡng ấy có cả
một dải, và trong dải đó mô hình dịch đúng gần như mọi câu ngắn rồi hỏng gần như
mọi câu dài.

\subsection*{Vì sao một seed một dòng là đủ ở đây và không đủ ở mục sau}

\index{đo lường!khi nào một lần rút là một xu hướng}%
Mục~\ref{sec:ch05-giai-ma} sẽ cho thấy ba mô hình chỉ khác nhau ở seed cho
0.8733, 0.6767 và 0.8700, tức tản gần hai mươi điểm. Với độ tản như thế, một
bảng chênh nhau vài điểm mà chỉ chạy một seed thì không đo được gì.

Bảng trên không ở tình huống đó. Hiệu ứng giữa hai đầu bảng là cả quãng từ 0 tới
0.8533, và mỗi bước tăng gấp đôi bề rộng lại cho một bước nhảy lớn hơn độ tản
seed. Còn mục~\ref{sec:ch05-dao-cau-nguon} đo một hiệu ứng cỡ mười tới ba mươi
điểm, nên nó chạy ba seed mỗi vế và in ra từng seed một. Chọn số lần lặp theo cỡ
hiệu ứng chứ không theo một quy tắc cố định, và nói ra mình chọn thế nào, là
phần rẻ nhất của việc đo đạc.

\subsection*{Cái bảng này chưa nói}

Nó không nói bài báo sai. Mô hình của bài mang 8000 số thực cho một câu, tức xa
bên phải mọi dòng trong bảng, và câu WMT dài hơn câu ở đây nhưng không dài hơn
theo tỷ lệ 8000 so với 256. Bảng này nói rằng chỗ thắt là có thật và đo được;
nó không nói mọi mô hình đều đã chạm vào nó.

Chuyện ai chạm vào nó, ai không, và ai là người đầu tiên viết ra điều đó, là
mục~\ref{sec:ch05-con-lai-gi}. Trước đó còn hai thứ bài làm mà tôi muốn kiểm,
và cả hai đều không đụng tới bề rộng vector.
